\section{Methods}

\subsection{Protocol and Registration}

The synthesis protocol was registered prospectively with the OpenTrials Registry
(registration OTR-2024-1187) prior to data extraction. Reporting follows the PRISMA 2020
guidance for systematic reviews and meta-analyses \citep{page2021prisma}.

\subsection{Search Strategy}

We searched the Global Research Index, ScholarNet, and the OpenTrials Registry for records
dated between January 2020 and December 2024, using combinations of the terms
\emph{``four-day workweek''}, \emph{``compressed schedule''}, \emph{``work-time
reduction''}, and \emph{``productivity pilot''}. Reference lists of retrieved records were
hand-searched for additional eligible trials. No language restriction was applied; all
identified records were in English.

\subsection{Eligibility Criteria}

Trials were eligible if they (i) compared an intervention group working a compressed
four-day, pay-neutral schedule against a concurrent control group working a standard
five-day schedule within the same organization; (ii) reported a quantitative,
standardized productivity or output measure for both arms; (iii) ran for a minimum
observation window of 10 weeks; and (iv) reported sufficient statistics (means, standard
deviations, and group sizes, or an equivalent effect size with its standard error) to
compute a standardized mean difference. Trials using only self-reported satisfaction as
the sole outcome, with no output-based measure, were excluded.

\subsection{Data Extraction and Risk-of-Bias Assessment}

Two reviewers independently extracted sample sizes, outcome definitions, and summary
statistics from each trial, resolving discrepancies by discussion. Each trial was appraised
with a five-domain risk-of-bias instrument covering allocation comparability, outcome
measurement, missing data, selective reporting, and confounding by concurrent
organizational change. All six included trials were rated as \emph{moderate} risk of bias,
principally due to non-random assignment of teams to schedule condition.

\subsection{Statistical Analysis}

For each trial we computed Hedges' $g$ \citep{hedges1981distribution}, a bias-corrected
standardized mean difference, together with its standard error. Trial-level effects were
pooled using a random-effects model with between-study variance ($\tau^2$) estimated by the
DerSimonian--Laird method \citep{dersimonian1986random}. Heterogeneity was quantified with
Cochran's $Q$, the $I^2$ statistic \citep{higgins2003measuring}, and a 95\% prediction
interval for the true effect in a new setting. Small-study effects were assessed with
Egger's regression test \citep{egger1997bias} and visual inspection of a funnel plot
(Appendix~\ref{app:funnel}). A leave-one-out sensitivity analysis was run to confirm that no
single trial drove the pooled estimate. All computations used MetaSynth v4.2 with a
two-sided significance threshold of $\alpha = 0.05$.
